\documentclass[11pt]{article}
\usepackage{mathblatt}
% gesetzt von werkzeuge/setzer.py aus dem Lernweg (L2-1 · L2-2 · L2-3 · L2-4 · L2-5 · L2-6 · L2-7) und den Bankzeilen; Muster T6B.
% Präambel des Setzers (werkzeuge/setzer.py), wortgleich aus dem Hand-Satz T6B (bau/proben/2026-10-09/kathete-berechnen), dort aus K4W.
\makeatletter
\setlength{\mbpfbe}{0mm}% keine Punkte (Bauregel 6.4)
% eingerückt (Bauregel 4.7, 6.6): \gEin Einzug, \gSchrift eine Stufe kleiner
\newlength{\gEin}\setlength{\gEin}{0pt}
\let\gSchrift\relax
\renewcommand{\pfkreuzzeile}[1]{\par\noindent\hangindent3.2em\hangafter1\hspace*{1.6em}$\square$\ #1\par}
\newcommand{\gkreuz}[1]{$\square$\ #1\hspace{2em}}
% Bauregel 6.7: links, was man ansieht, rechts, was man tut; Spaltengrenze überall an derselben Stelle
\newsavebox{\gbox}\newlength{\gLinks}
\newcommand{\gzwei}[2]{\par\smallskip\noindent\sbox{\gbox}{#1}%
  \setlength{\gLinks}{\dimexpr0.5\textwidth-\gEin-\mbpfnr\relax}%
  \ifdim\wd\gbox>\dimexpr\gLinks-4mm\relax
    \usebox{\gbox}\par\smallskip\noindent\begin{minipage}{\linewidth}#2\end{minipage}%
  \else
    \begin{minipage}[t]{\gLinks}\vspace{0pt}\usebox{\gbox}\end{minipage}%
    \begin{minipage}[t]{\dimexpr\linewidth-\gLinks\relax}\vspace{0pt}\raggedright #2\end{minipage}%
  \fi\par}
\RenewDocumentEnvironment{pfaufg}{m m m}{%
  \begin{lrbox}{\mbaufgabebox}\begin{minipage}[t]{\linewidth}%
  \gSchrift\noindent\hspace*{\gEin}\mbpfrandmarke{#1}%
  \begin{minipage}[t]{\mbpfnr}\strut\textbf{#2}\end{minipage}%
  \begin{minipage}[t]{\dimexpr\linewidth-\gEin-\mbpfnr\relax}\raggedright\strut\ignorespaces}%
 {\end{minipage}%
  \end{minipage}\end{lrbox}%
  \setlength{\mbaufgabehoehe}{\dimexpr\ht\mbaufgabebox+\dp\mbaufgabebox\relax}%
  \par\addvspace{7pt}\Needspace*{\mbaufgabehoehe}%
  \noindent\usebox{\mbaufgabebox}\par\addvspace{7pt}}
\makeatother
% Rechenraum als Karo (Bauregel 6.9): n Rechenschritte = n mal 10 mm
\newcommand{\karo}[1]{\par\smallskip\noindent\begin{tikzpicture}
  \clip (0,0) rectangle ({\linewidth-1mm},{#1*10mm});
  \draw[black!16,line width=0.3pt,step=5mm] (0,0) grid ({\linewidth},{#1*10mm});
  \end{tikzpicture}\par}
\newcommand{\kk}[1]{$\square$\,#1\hspace{0.9em}}
\newcommand{\linie}[1]{\,{\color{black!45}\rule[-1pt]{#1}{0.4pt}}\,}
% Zeile am Ende jedes Durchgangs: Originale klein grau, darüber; dann Vorgänger/Nachfolger (Bauregel 3.4, 6.5)
\newcommand{\blattende}[2]{\par\vfill\noindent\begin{minipage}{\linewidth}{\scriptsize\color{mbgrau}Originale: #1\par}\smallskip
{\footnotesize #2\par}\end{minipage}}
\newcommand{\pkt}{\textperiodcentered{}}

\newcommand{\kennung}{4NT}
\newcommand{\grau}[1]{{\color{mbgrau}#1}}

\begin{document}
\pfheftstil
\fancyfoot[C]{\parbox[t]{\textwidth}{\mbfhbox\par\vspace{1.5mm}{\footnotesize\color{mbgrau}{\scriptsize \kennung}\hfill\thepage}}}
\pfheftkopf{Prozentsatz berechnen}{Klasse 7}
% L2-1: Anschluss an Einheit 1: Der Streifen steht jetzt für eine Größe (40 kg); der Anteil ist Teil durch Ganzes, der Streifen bleibt 100 %
\begin{pfaufg}{}{1.}{}
Der ganze Streifen ist $100\,\%$. Wie viel Prozent sind grau? Schreibe zuerst den Bruch $\frac{\text{Teil}}{\text{Ganzes}}$.\par\medskip\noindent
\begin{minipage}[b]{0.32\linewidth}\centering
\begin{tikzpicture}
\fill[black!28] (0,0) rectangle (3.150,0.6);
\draw[black!55,line width=0.4pt] (1.050,0) -- (1.050,0.6);
\draw[black!55,line width=0.4pt] (2.100,0) -- (2.100,0.6);
\draw[black!55,line width=0.4pt] (3.150,0) -- (3.150,0.6);
\draw[line width=0.7pt] (0,0) rectangle (4.2,0.6);
\draw[line width=0.4pt] (0,0.6799999999999999) -- (0,0.82) -- (3.150,0.82) -- (3.150,0.6799999999999999);
\node[above,font=\small] at (1.575,0.8) {$30$ kg};
\node[below,font=\small] at (2.1,0) {ganz: $40$ kg};
\end{tikzpicture}\par\smallskip
\grau{{\small a)\ }\par\smallskip
$\frac{30}{40} = \frac{3}{4}$\par\smallskip $= 75\,\%$}
\end{minipage}\hfill
\begin{minipage}[b]{0.32\linewidth}\centering
\begin{tikzpicture}
\fill[black!28] (0,0) rectangle (1.680,0.6);
\draw[black!55,line width=0.4pt] (0.840,0) -- (0.840,0.6);
\draw[black!55,line width=0.4pt] (1.680,0) -- (1.680,0.6);
\draw[black!55,line width=0.4pt] (2.520,0) -- (2.520,0.6);
\draw[black!55,line width=0.4pt] (3.360,0) -- (3.360,0.6);
\draw[line width=0.7pt] (0,0) rectangle (4.2,0.6);
\draw[line width=0.4pt] (0,0.6799999999999999) -- (0,0.82) -- (1.680,0.82) -- (1.680,0.6799999999999999);
\node[above,font=\small] at (0.840,0.8) {$20\,€$};
\node[below,font=\small] at (2.1,0) {ganz: $50\,€$};
\end{tikzpicture}\par\smallskip
{\small b)\ }\par\smallskip
\linie{30mm}\par\medskip\linie{30mm}
\end{minipage}\hfill
\begin{minipage}[b]{0.32\linewidth}\centering
\begin{tikzpicture}
\fill[black!28] (0,0) rectangle (2.940,0.6);
\draw[line width=0.7pt] (0,0) rectangle (4.2,0.6);
\draw[line width=0.4pt] (0,0.6799999999999999) -- (0,0.82) -- (2.940,0.82) -- (2.940,0.6799999999999999);
\node[above,font=\small] at (1.470,0.8) {$140$ m};
\node[below,font=\small] at (2.1,0) {ganz: $200$ m};
\end{tikzpicture}\par\smallskip
{\small c)\ }\par\smallskip
\linie{30mm}\par\medskip\linie{30mm}
\end{minipage}
\fusshilfe{\mbox{1:~b) 40\,\%; c) 70\,\%}}
\end{pfaufg}

% L2-2: Erkennen ohne Rechnen: Was ist das Ganze? Es steht nicht immer im Text – mal ist es die Summe (13 + 3), mal der alte Preis, mal muss der Teil erst gebildet werden (frei = 24 − 18)
\begin{pfaufg}{}{2.}{}
Schreibe den Anteil als Bruch $\frac{\text{Teil}}{\text{Ganzes}}$. Rechne noch nicht aus.\par\medskip
\grau{a)\ $7$ von $28$ Kindern haben einen Hund. Anteil mit Hund?\quad $\frac{7}{28}$}\par\medskip
\noindent\begin{minipage}[t]{0.48\linewidth}
b)\ Von $600$ Losen sind $90$ Gewinne. Anteil der Gewinne?
\karo{1}
\end{minipage}\hfill
\begin{minipage}[t]{0.48\linewidth}
c)\ $13$ Pfannkuchen sind mit Zucker, $3$ mit Käse. Anteil mit Käse?
\karo{1}
\end{minipage}
\noindent\begin{minipage}[t]{0.48\linewidth}
d)\ Eine Jacke kostete $80\,€$, jetzt $60\,€$. Anteil des Rabatts am alten Preis?
\karo{1}
\end{minipage}\hfill
\begin{minipage}[t]{0.48\linewidth}
e)\ $45$ Gäste essen vegetarisch, $105$ nicht. Anteil vegetarisch?
\karo{1}
\end{minipage}
\noindent\begin{minipage}[t]{0.48\linewidth}
f)\ $18$ von $24$ Plätzen sind belegt. Anteil der freien Plätze?
\karo{1}
\end{minipage}
\fusshilfe{\mbox{2:~b) $\frac{90}{600}$; c) $\frac{3}{16}$; d) $\frac{20}{80}$; e) $\frac{45}{150}$; f) $\frac{6}{24}$}}
\end{pfaufg}

% L2-3: Ganzes ist ein Teiler oder Vielfaches von 100: wie in Einheit 1 erweitern (oder kürzen) bis Nenner 100
\begin{pfaufg}{}{3.}{}
Bringe den Bruch auf den Nenner $100$ (erweitern oder kürzen) und schreibe in Prozent.\par\medskip
\grau{a)\ $7$ von $20$\quad $\frac{7}{20} = \frac{35}{100} = 35\,\%$}\par\medskip
\noindent\begin{minipage}[t]{0.48\linewidth}
b)\ $19$ von $50$
\karo{1}
\end{minipage}\hfill
\begin{minipage}[t]{0.48\linewidth}
c)\ $9$ von $25$
\karo{1}
\end{minipage}
\noindent\begin{minipage}[t]{0.48\linewidth}
d)\ $3$ von $10$
\karo{1}
\end{minipage}\hfill
\begin{minipage}[t]{0.48\linewidth}
e)\ $84$ von $200$
\karo{1}
\end{minipage}
\noindent\begin{minipage}[t]{0.48\linewidth}
f)\ $18$ von $300$
\karo{1}
\end{minipage}
\fusshilfe{\mbox{3:~b) 38\,\%; c) 36\,\%; d) 30\,\%; e) 42\,\%; f) 6\,\%}}
\end{pfaufg}

% L2-4: Beliebiges Ganzes: Teil : Ganzes mit dem Taschenrechner, die Dezimalzahl als Prozent lesen; Kontrolle: Teil kleiner als Ganzes heißt unter 100 %
\begin{pfaufg}{}{4.}{}
Rechne mit dem Taschenrechner Teil : Ganzes und schreibe in Prozent.\par\medskip
\grau{a)\ $7$ von $20$\quad $\frac{7}{20} = 7 : 20 = 0{,}35 = \frac{35}{100} = 35\,\%$}\par\medskip
\noindent\begin{minipage}[t]{0.48\linewidth}
b)\ $18$ von $45$
\karo{1}
\end{minipage}\hfill
\begin{minipage}[t]{0.48\linewidth}
c)\ $51$ von $340$
\karo{1}
\end{minipage}
\noindent\begin{minipage}[t]{0.48\linewidth}
d)\ $117$ von $180$
\karo{1}
\end{minipage}\hfill
\begin{minipage}[t]{0.48\linewidth}
e)\ $9$ von $12$
\karo{1}
\end{minipage}
\noindent\begin{minipage}[t]{0.48\linewidth}
f)\ Ben rechnet bei $18$ von $45$ so: $45 : 18 = 2{,}5 = 250\,\%$. Woran siehst du ohne Rechnung, dass das falsch ist?
\karo{2}
\end{minipage}
\fusshilfe{\mbox{4:~b) 40\,\%; c) 15\,\%; d) 65\,\%; e) 75\,\%; f) unter 100\,\%}}
\end{pfaufg}

% L2-5: Das Ergebnis endet nicht: auf eine Stelle nach dem Komma runden; die dritte Stelle der Dezimalzahl entscheidet
\begin{pfaufg}{}{5.}{}
Rechne Teil : Ganzes und runde auf eine Stelle nach dem Komma.\par\medskip
\grau{a)\ $5$ von $7$\quad $5 : 7 = 0{,}7142\ldots = 71{,}42\ldots\,\% \approx 71{,}4\,\%$}\par\medskip
\noindent\begin{minipage}[t]{0.48\linewidth}
b)\ $13$ von $24$
\karo{1}
\end{minipage}\hfill
\begin{minipage}[t]{0.48\linewidth}
c)\ $41$ von $55$
\karo{1}
\end{minipage}
\noindent\begin{minipage}[t]{0.48\linewidth}
d)\ $11$ von $30$
\karo{1}
\end{minipage}\hfill
\begin{minipage}[t]{0.48\linewidth}
e)\ $7$ von $9$
\karo{1}
\end{minipage}
\fusshilfe{\mbox{5:~b) 54,2\,\%; c) 74,5\,\%; d) 36,7\,\%; e) 77,8\,\%}}
\end{pfaufg}

% L2-6: Das Ganze erst bilden: aus einer Tabelle (Summe aller Zeilen) und beim Rabatt (Differenz, bezogen auf den alten Preis)
\begin{pfaufg}{}{6.}{}
Alle Kinder der Klasse 7b haben gesagt, wie sie zur Schule kommen. Die Tabelle zeigt, wie viele Kinder es jeweils sind.\par\medskip\noindent \begin{minipage}[t]{0.40\linewidth}\vspace{0pt}{\renewcommand{\arraystretch}{1.3}\begin{tabular}{|l|c|}\hline
Schulweg & Kinder\\\hline Bus & $9$\\\hline Fahrrad & $12$\\\hline zu Fuß & $5$\\\hline Auto & $2$\\\hline\end{tabular}}\end{minipage}\hfill\begin{minipage}[t]{0.56\linewidth}\vspace{0pt}\raggedright Runde auf eine Stelle nach dem Komma.\par\medskip \textbf{a)}\ Wie viel Prozent der Klasse kommen mit dem Fahrrad?\linie{14mm}\,\%\par\medskip \textbf{b)}\ Wie viel Prozent kommen mit dem Auto?\linie{14mm}\,\%\end{minipage}\par
\karo{3}
\fusshilfe{\mbox{6:~a) 42,9\,\%; b) 7,1\,\%}}
\end{pfaufg}

% L2-6: Das Ganze erst bilden: aus einer Tabelle (Summe aller Zeilen) und beim Rabatt (Differenz, bezogen auf den alten Preis)
\begin{pfaufg}{}{7.}{}
Ein Fahrrad kostete $240\,€$. Im Angebot kostet es $192\,€$. Wie viel Prozent Rabatt gibt es?\par
\karo{3}
\fusshilfe{\mbox{7:~20\,\%}}
\end{pfaufg}

% L2-7: Zielaufgabe: Im Sachtext selbst sehen, ob Teil oder Ganzes erst gebildet werden muss, dann rechnen, runden und mit einem Satz antworten; das Ganze wechselt mit der Frage
\begin{pfaufg}{}{8.}{}
Die $30$ Kinder der 7b stimmen über das Ziel des Wandertags ab. $14$ Kinder sind für den Zoo, $10$ für den Kletterpark, $6$ Kinder stimmen nicht ab. Prüfe mit einer Rechnung, ob die Aussage stimmt.\par\smallskip
\noindent\begin{minipage}[t]{0.48\linewidth}
a)\ Lena sagt: „Mehr als die Hälfte der Klasse ist für den Zoo.“
\karo{3}
\end{minipage}\hfill
\begin{minipage}[t]{0.48\linewidth}
b)\ Tim sagt: „Von den Kindern, die abgestimmt haben, sind mehr als $40\,\%$ für den Kletterpark.“
\karo{3}
\end{minipage}
\fusshilfe{\mbox{8:~a) nein; b) ja}}
\end{pfaufg}

% L2-7: Zielaufgabe: Im Sachtext selbst sehen, ob Teil oder Ganzes erst gebildet werden muss, dann rechnen, runden und mit einem Satz antworten; das Ganze wechselt mit der Frage
\begin{pfaufg}{nach P10 ’15}{9.}{}
Mias Handballteam bekam in einer Saison $46$ Siebenmeter. $39$ davon wurden zu Toren. Berechne, wie viel Prozent der Siebenmeter nicht verwandelt wurden. Runde auf eine Stelle nach dem Komma.\par
\karo{3}
\fusshilfe{\mbox{9:~15,2\,\%}}
\end{pfaufg}

\par\vfill\noindent{\footnotesize Vorher: Prozente als Anteile\quad\textbullet\quad Weiter: Prozentwert berechnen\par}
\end{document}
